\section{Introduction}
\label{sec:introduction}

% A spacecraft approaching a small body initially sees a changing silhouette. Turning that silhouette into a reliable three-dimensional scientific model is essential to the scientific investigations and operations that follow. The dimensions of a depression, the slope beside a boulder, and the geometry of a narrow neck can influence both geological interpretation and the choice of where a spacecraft may safely operate.

For missions to explore distant small astroids(like NASA's OSIRIS-REx\citep{alasad2021}, ESA's Rosetta\citep{taylor2017rosetta}, JAXA's Hayabusa\citep{fujiwara2006,demura2006}), reliable three-dimensional models is essential throughout multiple scientific tasks. The process of obtaining models with captured photos by equipped optical instruments generates useful quantitative measurements: they support optical navigation, provide the spatial reference for mapping, place local observations within a global geological context, support scientific interpretation and sampling operations. Stereophotoclinometry (SPC)\citep{barnouin2020} is an established approach to extract information from captured obtical images and thus derive models of astroids. It connects pixels of multiple images and illumination status to stereo constraints like elevation and albedo of a specific location, which become integratable and optimizalbe variables.Its practical use across planetary missions demonstrates its value of extrating data in roughly unknow environments with limited informations.

Because the surface topography and albedo are totally or partially unknown before reconstruction, the observation geometry is the only controllable factor that can be used to improve the quality of reconstruction, therefore the quality of modeling depends on both the reconstruction procedure and the observations made available to it. The brightness of a pixel on images depens on not only the surface topography but also the observation geometry(illumination conditions and camera extrinsic parameters). Stereo observations of multiple images constrain the three-dimensional positions of the surface, while photoclinometry estimates local slopes and relative albedo from image brightness; these slopes are integrated into surface elevations with geometric and topographic constraints. Repeatedly imaging a specific area with the same illumination condition may therefore increase its nominal coverage without adequately constraining its elevation\citep{gaskell2023,palmer2022}, which makes observation of low infomation density. This limitation makes the choice of images a scientific determinant of model quality: an observation is valuable insofar as it adds constraints that the existing image set lacks. All above makes observation planning, the outcome of a sequence of operational decisions, an important scientific task for SPC model reconstruction.

SPC modeling proceeds through successive stages of refinement, starting from an initial global shape and progressively developing more detailed local terrain(called maplets). As images with closer spatial distance become available, the existing solution supports the construction of maplets with higher resolution, allowing finer surface structure to be established \citep{barnouin2020,palmer2022}. This progression changes the information needed from subsequent observations: early stages establish the overall shape and geometric framework, whereas later stages must strengthen local constraints and address deficiencies exposed by the preceding reconstruction. For a particular region of the current model, effective planning must answer three linked questions: \emph{Which observations are still missing, how much improvement could they provide, and how can that improvement be verified?} Existing SPC studies recommend complementary viewing directions and illumination conditions, including combinations of images that emphasize topography and relative albedo \citep{palmer2022,palmer2024geometry}. These recommendations provide valuable guidance for designing an observation plan, but a prescribed geometric configuration does not itself update as the reconstruction evolves or identify which additional observation would address the remaining deficiencies of a particular region. That decision requires the existing image set and the latest model-quality assessment to be considered together. Previous validation studies already provide spatial diagnostics, including maplet residuals and vertex height dispersion, that help identify regions with weak or inconsistent reconstruction support \citep{alasad2021,ernst2023}. Translating these diagnostics into the next acquisition remains a distinct planning problem: a region with uncertain heights may require a stronger stereo baseline, additional illumination diversity, or an observation that constrains albedo, and a large residual alone does not determine which intervention is appropriate. The cited studies do not establish a general feedback rule that makes this choice after each reconstruction round. A further requirement is an acquisition metric that predicts the scientific value of a candidate observation before its image is available. For the modeling objective considered here, that value can be expressed as the expected reduction in regional surface-height error or the increased likelihood of satisfying a prescribed accuracy tolerance, conditional on the current model and previously acquired images. Coverage and geometric diversity provide useful evidence, but their relationship to these accuracy gains requires validation. Such a metric would connect the diagnosis of existing deficiencies to the ranking of feasible future observations. After acquisition and reconstruction, the predicted gains must be checked against reference geometry in controlled experiments or independent measurements where available, with operational quality indicators calibrated accordingly \citep{weirich2022,alasad2021}. This connection between regional diagnosis, predicted observation benefit, and verified reconstruction improvement motivates the closed-loop planning approach studied here.



% Review version A: paragraph before "try again".
% The need for observation planning follows from how SPC recovers surface heights. 
 It also motivates studying observation sequences, since an intermediate reconstruction can reveal deficiencies that should change what is acquired next. A useful planning method must consequently connect candidate image combinations to the accuracy they are expected to support, then assess the resulting surface against independent evidence or validated quality indicators; agreement with the fitted images alone does not establish geometric accuracy \citep{weirich2022,alasad2021}. Establishing this connection is particularly necessary when target shape changes the available viewing geometry and when storage and communication limits restrict the images that can be acquired and delivered. These considerations motivate a reconstruction-driven planning framework that adapts across different small-body shapes, targets the unresolved needs of each modeling round, and allocates limited resources toward measurable accuracy requirements.

% Review version B: paragraph after "try again".
The reliability of an SPC model depends critically on whether the acquired images contain enough complementary information to resolve the surface geometry. A dark patch, for example, may reflect a surface tilted away from the Sun or material with a lower albedo. SPC addresses this ambiguity by combining stereo constraints on three-dimensional position with photometric estimates of local slope and relative albedo, then integrating the slopes into surface heights under geometric and topographic constraints \citep{gaskell2023,palmer2022}. Consequently, the recoverable terrain depends on the viewing directions and illumination conditions represented in the image set. Repeated observations with similar geometry can leave the same ambiguities unresolved, even as image count and surface coverage increase; sensitivity experiments demonstrate the importance of complementary observation geometries for SPC reconstruction \citep{palmer2024geometry}. This creates a practical scientific need to determine which image combinations constrain surface heights sufficiently to meet a specified accuracy requirement, and how reliably that accuracy can be anticipated before acquisition. The order of acquisition also matters in an adaptive campaign: observations obtained early determine which deficiencies can be diagnosed and addressed in subsequent rounds. Without a validated relationship between observation geometry and reconstruction error, a planner cannot reliably decide where further imaging is needed or when a region has been modeled accurately enough. Independent validation and calibrated quality indicators are therefore essential to this decision process \citep{weirich2022,alasad2021}. With finite storage and communication capacity, these considerations motivate automatic planning that uses reconstruction feedback to acquire complementary evidence and progressively establish a verifiable level of surface accuracy.

% Local topographic and relative-albedo maps are registered and assembled into larger terrain products, while estimates of observing geometry can also be refined \citep{gaskell2008,gaskell2023}.
%  However, the images supplied to SPC are themselves the outcome of a sequence of operational decisions. The quality of achievable model therefore depends on both the reconstruction procedure and the observations made available to it, whichmakes observation planning a scientific component of reconstruction. Consider a surface patch observed repeatedly from similar directions and under similar illumination. Its image count and nominal coverage can increase while important geometric ambiguities remain weakly constrained. A later observation that changes the relevant illumination or viewing geometry may provide information that the earlier images could not supply. SPC sensitivity experiments have demonstrated the importance of selecting complementary topographic and albedo observations, including a compact five-image configuration under the conditions tested for OSIRIS-REx \citep{palmer2024geometry}. Such results motivate a planner that reasons about the collective usefulness of an image set. They also raise a question that a coverage score alone cannot answer: how much geometric accuracy should be expected from the observations already acquired, and which additional observations are most likely to improve it?

The problem becomes more demanding as target morphology changes. SPC has already been applied to irregular bodies, including the bilobate nucleus of comet 67P/Churyumov--Gerasimenko \citep{jorda2016}. The generalization challenge considered here concerns the ability to use a consistent reconstruction and planning procedure across markedly different shapes, with limited target-specific adjustment. Elongation, deep concavities, steep terrain, and mutual occlusion alter which image combinations are attainable for each surface region. A coarse model may also misrepresent the very features that determine visibility and shadowing. Moreover, sensitivity tests identify steep slopes as a source of substantial local SPC errors \citep{palmer2024accuracy}. An observation policy must therefore respond to the target's geometry and the reconstruction's limitations, while recognizing regions that remain inaccessible within the available observing campaign.

An adaptive policy needs feedback that describes model quality in meaningful terms. After a reconstruction round, the planner should be able to identify regions with insufficient geometric support, inconsistent image evidence, or persistent uncertainty, and use those deficiencies to select the next observations. Establishing such feedback requires care. Internal agreement between overlapping terrain estimates is useful diagnostic evidence, but its relationship to absolute geometric accuracy must be assessed. Synthetic-asteroid tests explicitly investigated that relationship \citep{weirich2022}, while the validation of Bennu's SPC models combined several forms of quality assessment \citep{alasad2021}. Consequently, a per-round accuracy indicator should carry a clear spatial scale and an empirically tested interpretation. A small fitting residual or a high view count, by itself, cannot establish that the recovered surface is metrically correct.

The central scientific question of this work is therefore: \emph{Which observation sequences and image combinations enable SPC to produce high-accuracy small-body models whose accuracy can be predicted before acquisition and verified after reconstruction?} Here, predictability means estimating the quality attainable from a candidate image combination under stated geometric and measurement assumptions. Verifiability means checking the resulting geometric accuracy against independent evidence when available, and testing the reliability of operational quality indicators against reference geometry in controlled experiments. The sequence matters because it determines when new information becomes available for subsequent planning. The image combination matters because measurements of the same terrain can constrain different, complementary aspects of its shape and reflectance.

Existing research provides important components of a solution. Learned next-best-view methods estimate the value of candidate observations, while set-based approaches select multiple views together \citep{zeng2020,guedon2022,hu2024,pan2025}. Small-body imaging planners already account for photometric geometry and limited acquisition resources. In particular, \citet{piccinin2022} developed a reinforcement-learning policy for selecting imaging epochs, and \citet{wang2025} formulated observation planning for small-body photometric modeling. These studies establish that autonomous image selection and modeling-aware planning have precedents. The question pursued here is more specific: how to use the deficiencies of the \emph{current SPC reconstruction} to choose the next image set, and how to relate that choice to an accuracy assessment that remains meaningful across different body shapes.

Reinforcement learning offers a natural framework for the sequential part of this problem. An acquisition changes the available image combination, the knowledge of the surface, and the resources remaining for later observations. Its usefulness can therefore emerge over several decisions. For example, an image that contributes little new surface coverage may supply the missing viewing direction needed to constrain a terrain patch once combined with existing observations. Conversely, spending the remaining recorder capacity on immediately attractive images may prevent acquisition of a later, complementary view. A policy trained on the consequences of observation sequences can seek this longer-term scientific return. Scan-RL and GenNBV demonstrate the relevance of learned policies to active reconstruction and transfer across previously unseen geometries \citep{peralta2021,chen2024}; their results motivate investigating how sequential learning should be coupled to SPC-specific feedback.

The geometry of the action is equally important. A fixed catalogue of candidate viewpoints samples only part of the available observing configurations, and its angular spacing may be too coarse for a narrow visibility opening or a useful stereo baseline. In this work, we use a \emph{modified PointNet as a prediction model for continuous camera extrinsic parameters}, coupled with reinforcement learning for observation planning. The predicted extrinsics specify the position and orientation of a desired observation in a consistently defined coordinate frame. PointNet supplies a way to extract geometric features directly from an unordered point set \citep{qi2017pointnet}; its adaptation here makes continuous observation-pose prediction the network's task. Reinforcement learning supplies the sequential decision objective, relating the poses proposed through this model to the reconstruction achieved over a campaign.

This division of roles is central to the framework. The PointNet model predicts where and how to observe; SPC reconstructs the surface from the resulting image collection; the assessment of that reconstruction provides feedback for subsequent decisions. The pose output itself carries no automatic guarantee of terrain accuracy. Instead, the learning problem is organized around the relationship between the current reconstruction, its unresolved observation needs, and the eventual quality obtained after executing a sequence. Continuous pose prediction permits adaptation to changing geometry without fixing every possible observation in advance. Its practical benefit must still be established by comparing the resulting SPC models under equivalent acquisition opportunities and resource budgets.

Communication and storage constraints make this question operationally consequential. An image acquired today occupies recorder space and consumes future downlink capacity; its contribution to a ground-based reconstruction becomes available only after delivery and processing. Updated plans or model products may, in turn, require an uplink opportunity. Under a different allocation of computation between spacecraft and ground, the timing and content of these exchanges change, but their resource costs remain relevant. Mission-level autonomy studies emphasize the need to connect scientific intent, execution, and the reconstruction of onboard decisions \citep{castano2022}. For SPC, a planner must additionally consider whether a retained or transmitted image completes a useful geometric combination. Favorable acquisition geometry alone does not ensure that the data needed for the next reconstruction round will be available in time.

Our framework connects geometric representation, continuous observation-pose prediction, sequential decision making, and reconstruction feedback. At each planning round, the current surface representation provides the geometric context; available quality diagnostics and the observation history describe what is already constrained; and the mission state specifies the remaining opportunities and resources. The modified PointNet and reinforcement-learning components use this information within a repeated planning and reconstruction loop. Future images remain subject to acquisition and delivery, so the next update uses only evidence available at that stage of the campaign. In this formulation, the three research objectives are addressed through the following connected design choices:
\begin{enumerate}
    \item \textbf{Geometric adaptation across target morphology.} Use a point-set representation and continuous camera-extrinsic prediction to condition observations on the current body's shape, including elongation, concavities, and local relief. This provides a common interface between differently shaped surfaces and the observation policy, supporting a consistent SPC workflow with reduced reliance on manually specified viewing patterns. Generalization is assessed on held-out bodies and different coarse-model conditions; permutation invariance alone does not establish it.
    \item \textbf{Reconstruction-driven adaptation with per-round quality assessment.} Return the preceding SPC model and its diagnostic evidence to the planning loop, so that subsequent pose predictions respond to unresolved regions and missing geometric support. Reconstruction quality is evaluated separately from pose prediction, with region-level indicators calibrated against geometric error. The intended outcome is a sequence that progressively satisfies local accuracy requirements, with a record of which requirements remain unmet after each round.
    \item \textbf{Sequential allocation of communication and storage resources.} Couple the reinforcement-learning planning objective to uplink opportunities and capacity, downlink capacity, and recorder occupancy. This allows the scientific value of a pose to be considered together with the cost and timing of acquiring, retaining, transmitting, and using its image. Executable observations must satisfy the mission constraints; task completion is defined by the required reconstruction quality within those constraints, with infeasible requirements identified explicitly.
\end{enumerate}

These design choices explain how the proposed combination addresses the three research priorities, while also identifying what must be measured. A point-based pose predictor provides a geometry-dependent observation mechanism; feedback from successive SPC models makes the decisions responsive to actual reconstruction progress; and sequential resource accounting connects those decisions to task completion. Their scientific value rests on the reliability of regional accuracy assessments and their transfer across targets. The relevant evidence includes the surface area meeting a prescribed accuracy, errors in difficult terrain, the number of reconstruction rounds, image volume, and peak storage occupancy. Comparisons with fixed observing rules, discretized viewpoint policies, and continuous policies without reconstruction feedback can isolate the contribution of each component.

The resulting perspective is that an observation sequence should progressively remove the ambiguities that limit a scientific model. Automation becomes necessary as the surface representation, the useful viewing geometry, and the available resources evolve together. We study modified-PointNet pose prediction and reinforcement learning as a means to organize this coupled process, with SPC providing the reconstruction task and model validation providing the standard of success. The following review develops the connections to point-set learning, reinforcement-learning view planning, and spacecraft observation scheduling.
